% !TeX spellcheck = en_US
\documentclass[french]{yLectureNote}

\title{Algèbre linéaire 3*}
\subtitle{Physique}
\author{Paulhenry Saux}
\date{\today}
\yLanguage{Français}
%lombardi@math.univ-toulouse.fr
\professor{E.Lombardie}
\usepackage{graphicx}%----pour mettre des images
\usepackage[utf8]{inputenc}%---encodage
\usepackage{geometry}%---pour modifier les tailles et mettre a4paper
%\usepackage{awesomebox}%---pour les boites d'exercices, de pbq et de croquis ---d\'esactiv\'e pour les TP de PC
\usepackage{tikz}%---pour deiffner + d\'ependance de chemfig
% \usepackage{tabularx}%---pour dimensionner automatiquement les tableaux avec variable X
\usepackage{awesomebox}%---Pour les boites info, danger et autres
\usepackage{menukeys}%---Pour deiffner les touches de Calculatrice
\usepackage{fancyhdr}%---pour les en-t\^ete personnalis\'ees
\usepackage{blindtext}%---pour les liens
\usepackage{hyperref}%---pour les liens (\`a mettre en dernier)
\usepackage{caption}%---pour la francisation de la l\'egende table vers Tableau
\usepackage{pifont}
\usepackage{array}%---pour les tableaux
\usepackage{lipsum}
\usepackage{yFlatTable}
\usepackage{multicol}
\newcommand{\Lim}[1]{\lim\limits_{\substack{#1}}\:}
\renewcommand{\vec}{\overrightarrow}
\newcommand{\N}[0]{\mathbb{N}}
\newcommand{\R}[0]{\mathbb{R}}
\newcommand{\dd}{\mathrm{d}}
\newcommand{\norm}[1]{||\vec{#1}||}
\newcommand{\fo}{\psi(\vec{r},t)}
\newcommand{\foe}{\psi(\vec{r},t)\*}
\newcommand{\HH}{\hat{H}}
\newcommand{\hb}{\hbar}
\newcommand{\lap}{\nabla^2}
\newcommand{\lapcc}{\frac{\partial^2 }{\partial x^2}+\frac{\partial^2 }{\partial y^2}+\frac{\partial^2 }{\partial z^2}}
\newcommand{\E}{\vec{E}(M)}
\newcommand{\B}{\vec{B}(M)}
\newcommand{\V}{V(M)}
\newcommand{\er}{\vec{e_{\rho}}}
\newcommand{\ep}{\vec{e_{\varphi}}}
\newcommand{\pip}{\pi^+}
\newcommand{\pim}{\pi^-}
\newcommand{\mv}{\mathcal{V}}
\DeclareMathOperator{\Div}{Div}
\newcommand{\rot}{\vec{\mathrm{rot}}}
\newcommand{\grad}{\vec{\mathrm{grad}}}
\begin{document}
\chapter{Espace euclidien}
%C2 : Applications linéaires sur espaces euclidiens
%C3 : Espaces hermitiens
%C4 : Formes quadratiques
%C5 : Introduction aux groupes

%TD = U-111/U-112
%CM : B-08/B-07/Grignard

%CC2 : TD du 19 mars
%CC1 : 1h30
%CC2 : TD du 7 mai
%CC3 : 2 juin
%CC4 : 9 juin
\section{Produit scalaire et norme}
\begin{definition}[Produit scalaire]
    Soit $E$ un ev sur $\R$ de dimebsion finie ou non.
    $$\varphi : E\times E \to \R, (x,y)\to \varphi(x,y)$$
    \begin{itemize}
        \item $\varphi$ est une forme (l'espace d'arrivée est $\R$) bilinéaire si $x\to \varphi(x,y)$ est linéaire, de m\^eme pour y
        \item $\varphi$ est symétriuque $\iff \varphi(x,y) = \varphi(y,x)$
        \item $\forall x, \varphi(x,x) \geq0$
        \item $\varphi$ est définie positive $\iff \forall x, \varphi(x,x)\geq 0 \wedge \varphi(x,x) = 0 \Rightarrow x=0$
    \end{itemize}
    On dit que $\varphi$ est un produit scalaire $\iff \varphi$ est une forme bilinéaire symétrique définie positive.
\end{definition}
\warningInfo{Règles de clalcul}{On a 
\begin{itemize}
    \item $\varphi(\lambda x+x',y) = \lambda \varphi(x,y)+\varphi(x',y)$
    \item $\varphi(x,\lambda y+y') = \lambda \varphi(x,y)+\varphi((x,y'))$
\end{itemize}
}
\begin{definition}[Espace Euclidien]
    SI la dimension de $E$ est finie et $\varphi$ produit scalaire, alors la paire $(E,\varphi)$ est un espace euclidien.

\end{definition}
\begin{definition}[espace Préhilbertien]
        Si la dimension de $E$ est infinie, et $\varphi$ produit scalaire, alors la paire $(E,\varphi)$ est un espace préhilbertien.
\end{definition}
\checkInfo{Exemples}{
\begin{itemize}
    \item $E=\R^n : (x|y) = \sum_i^n x_iy_i$
    \item $E=\R^n : (x|y) = \sum_\lambda i^n x_iy_i, \lambda > 0$
    \item $E=\R_n[X] : (x|y) = \sum_{i=0}^{n} a_ib_i$
    \item $E = M_n[\R]$, avec $A = [A_{ij}], B=[B_{ij}] : \sum_{ij} a_{ij} b_{ij}, tr(^tAB)$
    \item $ E = C^0([0,1]; \R) : (x|y) = \int_0^1 fg\dd t$
\end{itemize}
}
\criticalInfo{Notation}{On note le produit scalaire $<x|y>$ ou $(x|y)$}
\begin{definition}[Norme]
    On note la norme associée au produit scalaire $\sqrt{(x|x)}$.
\end{definition}
\begin{proposition}[Inégalité de Cauchy Swartz]
    On munit E d'un produit scalaire, alors $\forall(x,y), (x|y) \leq ||x||||y||$
\end{proposition}
\begin{myproof}
    $\forall(x,y),\forall t\in \R$, calculons $(x+ty | x+ty)$.

    \begin{flalign*}
        (x+ty | x+ty) &\geq 0 \forall t\in \R\\
        &= (x|x) + (x|ty)+(ty|x)+(ty|ty)\\
        &= ||x||^2 + 2 t(x|y)+t^2||y||^2\\
        &= f(t)
    \end{flalign*}
    L'application f(t) est un polyn\^ome du second degré. On a donc $f(t) \geq 0$.

    Donc $\Delta \leq 0 \iff 4(x|y)^2- \leq 0 \iff 4(x|y)^2 \leq 4||x||^2$

    On a bien $$(x|y)^2 \leq ||x||^2$$
\end{myproof}
\begin{proposition}[Propriétés de la norme]
    \begin{itemize}
        \item $\forall x\in E, ||x||\geq 0$
        \item $||x|| = 0 \Rightarrow x=0$
        \item $||\lambda x|| = |\lambda|||x||$
        \item $||x+y|| \leq ||x||+||y||$
    \end{itemize}
\end{proposition}
\begin{myproof}[De $||x+y|| \leq ||x||+||y||$]
    \begin{flalign*}
        ||x+y||^2 &= (x+y|x+y)\\
        &= ||x||^2+2(x|y)+||y||^2\\
        &\leq ||x||^2+2||x||||y||+||y||^2\\
        &\leq (||x||+||y||)^2
    \end{flalign*}
\end{myproof}
\begin{proposition}[Calcul du produit scalaire]
    \begin{itemize}
        \item $(x|y) = \frac12 (||x+y||^2-||x||^2-||y||^2)$
        \item $(x|y) = \frac12(-||x-y||^2+||x||^2+||y||^2)$
        \item $(x|y) = \frac14(||x+y||^2-||x-y||^2)$
        \item $||x+y||^2+||x-y||^2 = 2(||x||^2+||y||^2)$
    \end{itemize}
\end{proposition}
\begin{myproof}[Formules 2, 3 et 4]
    \textbf{2.} En remplaçant $y$ par $-y$ dans la formule 1 :
    \begin{flalign*}
        (x|-y) &= \frac12(||x-y||^2-||x||^2-||-y||^2)\\
        -(x|y) &= \frac12(||x-y||^2-||x||^2-||y||^2)\\
        (x|y) &= \frac12(-||x-y||^2+||x||^2+||y||^2)
    \end{flalign*}
    
    \textbf{3.} En additionnant les formules 1 et 2 :
    \begin{flalign*}
        2(x|y) &= \frac12(||x+y||^2-||x||^2-||y||^2) + \frac12(-||x-y||^2+||x||^2+||y||^2)\\
        2(x|y) &= \frac12(||x+y||^2-||x-y||^2)\\
        (x|y) &= \frac14(||x+y||^2-||x-y||^2)
    \end{flalign*}
    
    \textbf{4.} En additionnant les formules 1 et 2 sans la division par 2 :
    \begin{flalign*}
        ||x+y||^2 &= ||x||^2+2(x|y)+||y||^2\\
        ||x-y||^2 &= ||x||^2-2(x|y)+||y||^2\\
        ||x+y||^2+||x-y||^2 &= 2||x||^2+2||y||^2 = 2(||x||^2+||y||^2)
    \end{flalign*}
\end{myproof}
\section{Orthogonalité}
\subsection{Introduction}
\begin{definition}[Vecteurs orthogonaux]
    Deux vecteurs y et x sont orthogonaux quand leur produit scalaire est nul. On note $x\perp y$
\end{definition}
\begin{theorem}[Pythagore]
    En réels, $x\perp y\iff ||x+y||^2 = ||x||^2+||y||^2$
\end{theorem}
\begin{myproof}
    On calcule la différence :
    \begin{flalign*}
        ||x+y||^2-||x||^2-||y||^2 &= ||x||^2+||y||^2+2(x|y)-||x||^2-||y||^2\\
        &= 2(x|y)\\
        &=0
    \end{flalign*}
\end{myproof}
\begin{definition}
    Une base $e_1,\dots,e_n$ est orthogonale quand $\forall i,j, i\neq j, e_i\perp e_j$ 

    Une base est orthonormée quand $\forall i\neq j, e_i\perp e_j, \forall i, ||e_i||=1$
\end{definition}
On peut utiliser le symbole de $\delta_{ij}$
\begin{proposition}
    Si $e_1\dots e_n$ est un BON, alors $\forall x\in E, x= (e_1|x)e_1+\dots + (e_n|x)e_n$
    Les coordonnées de x dans la base $e_1,\dots, e_n$, notées $(x_1,\dots, x_n)$ sont données par $x_i = (e_i|x)$
\end{proposition}
\begin{corollary}
    $||x||^2 = x_1^2+\dots+x_n^2 = (e_1|x)^2+\dots+(e_n|x)^2$
\end{corollary}
\begin{myproof}
    \begin{flalign*}
        ||x||^2 &= (x|x)\\
        &= (x_1e_1+\dots|x_1e_1+\dots)\\
        &= \sum_{i,j}(x_i e_i|x_je_j)\\
        &= \sum_{i,j}x_ix_j(e_i|e_j)\\
        &= \sum_{i,j} x_ix_j\delta_{ij}\\
        &= \sum_{i} x_i^2\times 1\\
        &= \sum x_i^2
    \end{flalign*}
\end{myproof}
\subsection{Procédés d'orthonormalité de Gram-Schimdt}
\begin{proposition}
    Soit $(u_1, _dots, u_n)$ une base quelconque de $E$. Il existe une base orthonormée $e_1,\dots, e_n$
    telle que $\forall p, Vect(e_1,\dots, e_p) = Vect(u_1,\dots, u_p)$
\end{proposition}
\begin{myproof}[Par recurrence]
    Initialisation : $p=1$, on a $e_1 = \frac{u_1}{||u_1||}$

    Hérédité : Supposons que l'on ait construit $(e_1, \dots,e_p)$ orthornomée telle que 
    $Vect(e_1,\dots,e_p) = Vect(u_1,\dots, u_p)$.

    Construisons $e_{p+1}$. On veut $e_1,\dots,e_p$ BON et $Vect(e_1,\dots,e_{p+1}) = Vect(u_1,\dots, u_{p+1})$.

    On prend : 
    \begin{enumerate}
        \item $e'_{p+1} = u_{p+1}+\sum_{i=1}^p \alpha_i e_i$

    \item On veut que $e'_{p+1}\perp e_j \forall 1\leq j\leq p$

    \item On a $e_{p+1} = \frac{e'_{p+1}}{||e'_{p+1}||}$.
    \end{enumerate}
    On a 
    \begin{flalign*}
        0 &= (e'_{p+1}|e_j)\\
        \iff 0 &= (u_{p+1}+\sum_{i=1}^n \alpha_i e_i|e_j)\\
        \iff 0 &= (u_{p+1}|e_j) +(\sum_{i=1}^p \alpha_i e_i|e_j)\\
        \iff 0 &= (u_{p+1}|e_j) +\alpha_j\\
        \alpha_j = -(u_{p+1}|e_j)
    \end{flalign*}
\end{myproof}
\warningInfo{Méthode de Gram}{
    \begin{enumerate}
        \item $e_1 = \frac{u_1}{||u_1||}$
        \item $1\leq p\leq n-1, e'_{p+1} = u_{p+1} - \sum_{j=1}^p(e_j|u_{p+1})e_j$
        \item $e_{p+1} = \frac{e'_{p+1}}{||e'_{p+1}||}$
    \end{enumerate}
}
\section{Sous espace orthogonal et projection}
\begin{definition}
    Soit E un espace euclidien. A est une partie de E. L'orthogonal de A, noté 
    $A^{\perp}$ est donné par $A^{\perp} = \{y\in E/ \forall a \in A (y|a)=0\}$
\end{definition}
\begin{proposition}
    $A^{\perp}$ est toujours un sev de E m\^eme si A n'est qu'une partie quelconque de E.
\end{proposition}
\begin{myproof}
    $0\in A^{\perp}$  donc non vide.

    Soit $y_1, y_2\in A^{\perp}, \lambda\in \R$. On veut montrer que $\lambda y_1+y_2\in A^{\perp}$ 

    $y_1\in A^{\perp}\iff \forall a \in A, (y_1|a) = 0$

    $y_2\in A^{\perp}\iff \forall a \in A, (y_2|a) = 0$

    Donc $\forall a \in A, \lambda (\lambda y_1+y_2 | a) = 0$ par combinaison linéaire.
\end{myproof}
\begin{proposition}
    $F$ sev de E euclicien, F et $F^{\perp}$ sont supplémentaires $\iff E = F\oplus F^{\perp}$
\end{proposition}
\checkInfo{Somme directe}{
F et G sont en somme directe $\iff F\cap G = \{0\} \iff \forall z\in F+G, \exists! f\in F, g\in G, z=f+g$

F et G supplémentaires $\iff$ F et G en somme directe et $E = F+G$
}
\begin{myproof}
    \begin{enumerate}
        \item Montrons F et $F^{\perp}$ en somme directe.
        
        Soit $x\in F\cap F^{\perp}$. $||x||^2 = (x|x) = 0$ donc $x=0$

        Donc $F\cap F^{\perp}$ est en somme directe
        \item Montrons que $E=F+F^{\perp}$. Il faut montrer que $\forall x\in E, x_F\in F, x_{F^{\perp}} \in F^{\perp}, x= x_F  + x_{F^{\perp}}$
        
        On note p la dimension de F. Avec Gram-Schmit, on peut construire une base orthonormée $e_1,\dots, e_p$
        de F et comme c'est une bon, $x_F = \alpha_1e_1+\dots+\alpha_pe_p$. 

        Donc $1\leq j\leq p$, $(x|e_j) = (x_F+x_{\perp}|e_j) = \alpha_j$

        Donc nécessairement, $x_F = (e_1|x)e_1+\dots+(e_p|x)e_p$.

        Reste à vérifier que $x_F = (e_1|x)e_1+\dots+(e_p|x)e_p$ et $x_{\perp} = x-x_F$.
        \item Il ne reste plus qu'à démontrer que $x_{\perp} = x-x_F \in F^{\perp}$
        C'est équivalent à $\forall y\in F, (x_{\perp}|y)=0$

        On peut juster vérifier que c'est orthogonal à une base de F, i.e. $\forall j, (x_{\perp}|e_j)$.

        $(x_{\perp} |e_j) = (x-x_F|e_j) = (x|e_j)-(\sum (e_i|x)e_i|e_j) = 0$

        Donc $\forall j\in x_{\perp}\perp e_j$

        Donc OK
    \end{enumerate}
\end{myproof}
\begin{corollary}
    La projection orthogonale sur F est donnée par $p_F(x) = (e_1|x)e_1+\dots+(e_p|x)e_p$
    or $(e_1,\dots e_p)$ BON
\end{corollary}
\begin{corollary}
    $\dim(F) + \dim(F^{\perp})=\dim(E)$
\end{corollary}
\begin{corollary}
    $F^{\perp\perp} = F$
\end{corollary}
\begin{myproof}
    $F\subset F^{\perp \perp}$

    $\dim(F) + \dim(F^{\perp}) = \dim(E)$ et $\dim(F^{\perp}+\dim(F^{\perp\perp}))$. on soustrait, cela fait 0.

    Les dimensions sont finies, donc l'égalité est vérifie.
\end{myproof}
\begin{proposition}
    F un SEV de E, $x\in E$. Alors $\forall y \in F, ||x-p_f(x)||\leq||x-y|| \iff ||x-p_f(x)|| = \min(||x-y||)$
    C'est la distance de x à F et $p_f(x)$ est la projection orthogonale de x sur F.
\end{proposition}
\begin{myproof}
    $\forall x\in E, x = x_F+x_\perp = p_f(x)+h$

    De plus, $\forall y\in F, ||x-y||^2 = ||p_f(x)-y+h||^2 = ||y-p_f(x)||^2 + ||h||^2 = ||p_f(x)-y||^2+||x-p_f(x)||^2$
    Donc $\forall y, ||x-y||^2 = ||p_f(x)-y||+||x-p_f(x)||^2$

    Donc $||x-y||^2\geq ||x-p_f(x)||^2$. En prenant $y = p_f(x)\in F$, le minimum est atteint et la
    distance de x à F = $\min(||x-y||) = ||x-p_f(x)||$.
\end{myproof}

\end{document}
